\documentclass[11pt]{article}
\usepackage{amsmath,amssymb,amsthm}
\usepackage[margin=1in]{geometry}

\newtheorem{theorem}{Theorem}
\newtheorem{lemma}[theorem]{Lemma}

\title{Problem 642: Sum of Largest Prime Factors}
\author{Project Euler Solutions}
\date{}

\begin{document}
\maketitle

\section{Problem Statement}

Let $\operatorname{lpf}(n)$ denote the largest prime factor of~$n$ (with $\operatorname{lpf}(1)=1$). Compute
\[
S(N) = \sum_{n=2}^{N}\operatorname{lpf}(n).
\]

\section{Mathematical Foundation}

\begin{theorem}[Sieve Correctness]
If primes are processed in increasing order and for each prime~$p$ we set $\operatorname{lpf}[m]\leftarrow p$ for every multiple~$m$ of~$p$, then at termination $\operatorname{lpf}[m]$ equals the largest prime factor of~$m$ for every $2\le m\le N$.
\end{theorem}

\begin{proof}
Let $q=\operatorname{lpf}(m)$ be the true largest prime factor. Since $q\mid m$ and $q\le N$, the iteration for~$q$ sets $\operatorname{lpf}[m]\leftarrow q$. No prime $p>q$ divides~$m$, so no later iteration overwrites. Hence $\operatorname{lpf}[m]=q$.
\end{proof}

\begin{lemma}[Regrouping]
\[
S(N) = \sum_{\substack{p\le N\\p\text{ prime}}} p\cdot\#\{n\le N:\operatorname{lpf}(n)=p\}.
\]
\end{lemma}

\begin{proof}
Rearrange the sum by grouping integers sharing the same largest prime factor.
\end{proof}

\begin{theorem}[Sub-linear Approach]
Define $F(n,p)$ as the count of integers $\le n$ with all prime factors $\le p$. Then
\[
\#\{n\le N:\operatorname{lpf}(n)=p\} = F(N/p,\,p) - F(N/p,\,p^{-}),
\]
where $p^{-}$ is the prime preceding~$p$. This yields an $O(N^{2/3})$ algorithm via the Lucy~DP technique.
\end{theorem}

\begin{proof}
An integer $n\le N$ satisfies $\operatorname{lpf}(n)=p$ iff $p\mid n$ and $n/p$ has all prime factors $\le p$. The count of such $n/p\le N/p$ with largest prime factor $\le p$ is $F(N/p,p)$. Subtracting $F(N/p,p^{-})$ excludes those with largest prime factor $<p$.
\end{proof}

\begin{lemma}[Asymptotic]
$S(N)\sim N^2/(2\ln N)$.
\end{lemma}

\begin{proof}
By partial summation and the prime number theorem, the dominant contribution comes from integers whose largest prime factor is $\Theta(N)$. The Dickman function analysis gives the stated asymptotic.
\end{proof}

\section{Algorithm}

\subsection{Sieve Method}
\begin{enumerate}
\item Initialise $\operatorname{lpf}[0\ldots N]=0$.
\item For each prime~$p$ from~2 to~$N$: set $\operatorname{lpf}[m]=p$ for all multiples $m$ of~$p$.
\item Return $\sum_{n=2}^{N}\operatorname{lpf}[n]$.
\end{enumerate}

\subsection{Sub-linear Method}
Use the Lucy~DP framework to compute $F(v,p)$ for $v\in\{\lfloor N/k\rfloor:k=1,\ldots,N\}$ and apply the regrouping identity.

\section{Complexity Analysis}

\begin{itemize}
\item \textbf{Sieve:} Time $O(N\log\log N)$, Space $O(N)$.
\item \textbf{Sub-linear (Lucy DP):} Time $O(N^{2/3})$, Space $O(N^{1/2})$.
\end{itemize}

\section{Answer}

\[
\boxed{631499044}
\]

\end{document}
